\documentclass[11pt,a4paper]{article}

\usepackage[T1]{fontenc}
\usepackage[utf8]{inputenc}
\usepackage[brazilian]{babel}
\usepackage{amsmath,amssymb,mathtools}
\usepackage[protrusion=true,expansion=false]{microtype}
\usepackage{needspace}
\usepackage{hyperref}
\usepackage{bookmark}
\usepackage[most]{tcolorbox}
\usepackage{geometry}
\usepackage{enumitem}
\usepackage{xcolor}
\geometry{margin=2.2cm}
\hypersetup{colorlinks=true,linkcolor=blue!60!black,urlcolor=blue!60!black}
\setlist[itemize]{topsep=4pt,itemsep=2pt,leftmargin=1.5em}
\setlist[enumerate]{topsep=4pt,itemsep=2pt,leftmargin=1.7em}
\setlength{\parskip}{0.35em}
\setlength{\parindent}{0pt}
\setlength{\abovedisplayskip}{8pt plus 2pt minus 3pt}
\setlength{\belowdisplayskip}{8pt plus 2pt minus 3pt}
\setlength{\abovedisplayshortskip}{6pt plus 2pt minus 2pt}
\setlength{\belowdisplayshortskip}{6pt plus 2pt minus 2pt}
\clubpenalty=10000
\widowpenalty=10000
\displaywidowpenalty=10000
\brokenpenalty=10000
\raggedbottom
\newcommand{\dd}{\mathrm{d}}
\let\oldsection\section
\RenewDocumentCommand{\section}{s o m}{%
  \Needspace{12\baselineskip}%
  \IfBooleanTF{#1}
    {\oldsection*{#3}}
    {\IfNoValueTF{#2}{\oldsection{#3}}{\oldsection[#2]{#3}}}%
}
\let\oldsubsection\subsection
\RenewDocumentCommand{\subsection}{s o m}{%
  \Needspace{9\baselineskip}%
  \IfBooleanTF{#1}
    {\oldsubsection*{#3}}
    {\IfNoValueTF{#2}{\oldsubsection{#3}}{\oldsubsection[#2]{#3}}}%
}

\newtcolorbox{resultado}{
  colback=blue!4!white,
  colframe=blue!55!black,
  title=Resultado final,
  before=\par\Needspace{9\baselineskip}\medskip,
  after=\medskip,
  boxrule=0.7pt,
  left=7pt,
  right=7pt,
  top=5pt,
  bottom=5pt,
  breakable
}

\newtcolorbox{armadilha}{
  colback=red!4!white,
  colframe=red!55!black,
  title=Erro comum,
  before=\par\Needspace{8\baselineskip}\medskip,
  after=\medskip,
  boxrule=0.7pt,
  left=7pt,
  right=7pt,
  top=5pt,
  bottom=5pt,
  breakable
}

\newtcolorbox{roteiro}{
  colback=orange!5!white,
  colframe=orange!65!black,
  title=Roteiro de estudo,
  before=\par\Needspace{10\baselineskip}\medskip,
  after=\medskip,
  boxrule=0.7pt,
  left=7pt,
  right=7pt,
  top=5pt,
  bottom=5pt,
  breakable
}

\title{Guia de deriva\c{c}\~oes --- Cap\'itulos 1 e 2}
\author{romer-study}
\date{\today}

\begin{document}

\maketitle
\tableofcontents
\bigskip

\section{Objetivo deste guia}

Este PDF n\~ao substitui as notas de cada cap\'itulo. Ele funciona como um
guia de apoio para estudar as deriva\c{c}\~oes com menos atrito. A ideia \'e
responder a tr\^es perguntas:
\begin{enumerate}
  \item o que sempre precisa ser definido antes da conta come\c{c}ar;
  \item quais passos alg\'ebricos se repetem de Solow para RCK e Diamond;
  \item onde os alunos costumam se perder.
\end{enumerate}

\begin{roteiro}
Use este guia assim:
\begin{enumerate}
  \item leia primeiro a nota completa do cap\'itulo;
  \item volte aqui para revisar o \emph{esqueleto} da deriva\c{c}\~ao;
  \item refa\c{c}a as contas no papel sem olhar os passos intermedi\'arios;
  \item por fim, compare com o c\'odigo do reposit\'orio.
\end{enumerate}
\end{roteiro}

\section{Mapa de nota\c{c}\~ao}

\subsection{Tr\^es n\'iveis de vari\'aveis}

Ao longo dos Cap\'itulos 1 e 2, quase toda a confus\~ao vem de misturar tr\^es
tipos de vari\'aveis:

\begin{itemize}
  \item \textbf{agregadas}: \(Y\), \(K\), \(C\), \(L\), \(AL\);
  \item \textbf{por trabalhador}: \(Y/L\), \(C/L\);
  \item \textbf{por trabalhador efetivo}: \(y=Y/(AL)\), \(c=C/(AL)\), \(k=K/(AL)\).
\end{itemize}

\begin{resultado}
Antes de derivar qualquer equa\c{c}\~ao, pergunte:
\begin{quote}
``Esta vari\'avel est\'a em n\'ivel agregado, por trabalhador ou por trabalhador
efetivo?''
\end{quote}
\end{resultado}

\subsection{Correspond\^encias com o projeto}

\begin{itemize}
  \item Solow:
  \begin{itemize}
    \item \(f(k)\) \(\leftrightarrow\) \texttt{SolowModel.f(...)}
    \item \(f'(k)\) \(\leftrightarrow\) \texttt{SolowModel.f\_prime(...)}
    \item \(\dot{k}\) \(\leftrightarrow\) \texttt{SolowModel.k\_dot(...)}
  \end{itemize}
  \item RCK:
  \begin{itemize}
    \item \(\dot{k}\) \(\leftrightarrow\) \texttt{RCKModel.k\_dot(...)}
    \item \(\dot{c}\) \(\leftrightarrow\) \texttt{RCKModel.c\_dot(...)}
    \item saddle path \(\leftrightarrow\) \texttt{find\_saddle\_path(...)}
  \end{itemize}
\end{itemize}

\section{Template 1: como reduzir um modelo agregado}

\subsection{Passo geral}

Quando o modelo come\c{c}a com
\[
Y=F(K,AL),
\]
o roteiro quase sempre \'e este:
\begin{enumerate}
  \item usar retornos constantes de escala;
  \item definir \(k=K/(AL)\) e \(y=Y/(AL)\);
  \item escrever \(Y=ALf(k)\);
  \item derivar \(k\) no tempo com a regra do quociente.
\end{enumerate}

\subsection{A conta que mais se repete}

Se
\[
k=\frac{K}{AL},
\]
ent\~ao
\[
\dot{k}
=
\frac{\dot K}{AL}
-
\frac{K}{(AL)^2}\frac{\dd(AL)}{\dd t}.
\]
Como
\[
\frac{\dd(AL)}{\dd t}=AL(n+g),
\]
obtemos
\[
\dot{k}=\frac{\dot K}{AL}-(n+g)k.
\]

\begin{resultado}
Toda vez que voc\^e normaliza por \(AL\), aparecem automaticamente os termos
de ``dilui\c{c}\~ao'' \(nk\) e \(gk\).
\end{resultado}

\begin{armadilha}
Erro comum: esquecer que \(\dot{(AL)}/(AL)=n+g\), e n\~ao apenas \(g\) ou apenas \(n\).
\end{armadilha}

\section{Template 2: como obter a equa\c{c}\~ao de Euler}

\subsection{Roteiro do RCK}

No RCK, a sequ\^encia-padr\~ao \'e:
\begin{enumerate}
  \item escrever a restri\c{c}\~ao em \(k\) e \(c\);
  \item normalizar a utilidade;
  \item montar o Hamiltoniano;
  \item tirar a FOC em \(c\);
  \item usar a equa\c{c}\~ao do coestado;
  \item eliminar \(\lambda\).
\end{enumerate}

\subsection{Conta essencial}

Da FOC,
\[
\lambda=c^{-\theta}.
\]
Logo,
\[
\frac{\dot\lambda}{\lambda}=-\theta\frac{\dot c}{c}.
\]
Se a equa\c{c}\~ao do coestado entrega
\[
\frac{\dot\lambda}{\lambda}=\rho+\theta g+\delta-f'(k),
\]
ent\~ao basta igualar:
\[
-\theta\frac{\dot c}{c}
=
\rho+\theta g+\delta-f'(k).
\]
Portanto,
\[
\frac{\dot c}{c}
=
\frac{f'(k)-\delta-\rho-\theta g}{\theta}.
\]

\begin{armadilha}
Erro comum: esquecer que, na conven\c{c}\~ao por trabalhador efetivo, aparece o
termo \(\theta g\) na Euler equation.
\end{armadilha}

\section{Template 3: como transformar escolha de poupan\c{c}a em mapa din\^amico}

\subsection{Roteiro do Diamond}

No Diamond simples, a sequ\^encia \'e:
\begin{enumerate}
  \item escrever as duas restri\c{c}\~oes or\c{c}ament\'arias do jovem e do velho;
  \item substituir na utilidade;
  \item derivar a FOC em \(s_t\);
  \item escrever o sal\'ario como fun\c{c}\~ao de \(k_t\);
  \item dividir pela nova coorte para obter \(k_{t+1}\).
\end{enumerate}

\subsection{Conta essencial}

Se
\[
s_t=\frac{\beta}{1+\beta}w_t
\]
e
\[
w_t=(1-\alpha)k_t^\alpha,
\]
ent\~ao
\[
k_{t+1}
=
\frac{s_t}{1+n}
=
\frac{\beta(1-\alpha)}{(1+\beta)(1+n)}k_t^\alpha.
\]

\begin{resultado}
O Diamond simples sempre termina em um mapa discreto da forma
\[
k_{t+1}=\Phi(k_t),
\]
e n\~ao em uma equa\c{c}\~ao diferencial.
\end{resultado}

\begin{armadilha}
Erro comum: usar a mesma intui\c{c}\~ao do Solow ou do RCK para o Diamond. No
OLG, o problema \'e discreto e a demografia entra diretamente na passagem de
poupan\c{c}a agregada para capital por jovem.
\end{armadilha}

\section{Checklist de boas pr\'aticas ao derivar}

\begin{enumerate}
  \item Declare as vari\'aveis e a unidade de medida antes de come\c{c}ar.
  \item Escreva a restri\c{c}\~ao din\^amica antes da condi\c{c}\~ao de otimalidade.
  \item N\~ao pule a etapa de dividir por \(AL\) ou por \(L\).
  \item Quando aparecer um steady state, imponha explicitamente \(\dot{k}=0\),
  \(\dot{c}=0\) ou \(k_{t+1}=k_t\).
  \item Ao usar Cobb--Douglas, isole primeiro a pot\^encia de \(k\) e s\'o depois
  resolva a express\~ao fechada.
  \item Separe sempre o que \'e \textbf{conta alg\'ebrica} do que \'e
  \textbf{interpreta\c{c}\~ao econ\^omica}.
\end{enumerate}

\section{Sequ\^encia sugerida de estudo}

\begin{enumerate}
  \item Solow:
  \begin{itemize}
    \item derive \(y=f(k)\);
    \item derive \(\dot{k}\);
    \item resolva \(k^\ast\) e a Regra de Ouro;
    \item s\'o depois olhe o growth accounting.
  \end{itemize}
  \item RCK:
  \begin{itemize}
    \item derive a restri\c{c}\~ao normalizada;
    \item monte o Hamiltoniano;
    \item derive a Euler equation;
    \item feche com o diagrama de fase.
  \end{itemize}
  \item Diamond:
  \begin{itemize}
    \item derive a regra de poupan\c{c}a;
    \item derive o mapa \(k_{t+1}(k_t)\);
    \item compare \(k^\ast\) com \(k_{\text{gold}}\).
  \end{itemize}
\end{enumerate}

\section{Onde continuar no reposit\'orio}

Depois de ler este guia, o caminho natural \'e abrir:
\begin{itemize}
  \item \texttt{ch01\_solow/notes/ch01\_solow\_derivations.pdf};
  \item \texttt{ch02\_rck\_diamond/notes/ch02\_rck\_diamond\_derivations.pdf};
  \item \texttt{ch01\_solow/ch01\_solow.py};
  \item \texttt{ch02\_rck\_diamond/ch02\_rck.py}.
\end{itemize}

\end{document}
